\chapter{KẾT LUẬN VÀ HƯỚNG PHÁT TRIỂN}

\section{Kết quả giai đoạn giữa kỳ}

Nhóm đã hoàn thành khảo sát ba hướng mô hình, chuẩn bị dữ liệu đầu vào từ chú thích gốc và thực nghiệm chín cặp mô hình--bộ dữ liệu. BROS đạt F1 cao nhất trong các lượt chạy hiện có; RoBERTa là mốc so sánh văn bản có kết quả cạnh tranh trên SROIE; PICK cần được rà soát thêm trước khi diễn giải nguyên nhân điểm thấp. Kết quả được lưu cùng checkpoint và dự đoán chi tiết để có thể kiểm tra.

Báo cáo đã phân biệt bài báo tham khảo với bản triển khai, mô tả rõ metric đang dùng và phạm vi đánh giá thực nghiệm. Những giới hạn về một seed, cắt token và khác biệt của PICK được ghi nhận để tránh kết luận vượt quá bằng chứng.

\section{Công việc tiếp theo}

Ưu tiên tiếp theo là hoàn tất kiểm tra bản triển khai PICK và thống nhất cách xử lý tài liệu dài trước khi mở rộng thực nghiệm. Mỗi thay đổi cần được đánh giá bằng validation và ghi rõ cấu hình; không dùng lỗi trên test để lựa chọn siêu tham số. Nếu tiếp tục chạy thêm seed, nhóm sẽ báo cáo trung bình và mức biến động thay vì chỉ một giá trị.

Nhóm sẽ hoàn thiện phân tích lỗi có ảnh minh họa, bổ sung thời gian chạy và chi phí tài nguyên khi có phép đo đầy đủ, đồng thời kiểm tra lại cách chấm tham chiếu của từng bộ. Các kết quả bổ sung chỉ đưa vào bản báo cáo sau khi đã thực sự chạy và đối chiếu. Phạm vi đề tài vẫn là đánh giá và so sánh phương pháp trích xuất thông tin từ hóa đơn.
